% リアクションホイールと手足
\begin{tikzpicture}[>=Stealth,font=\small]
\begin{scope}
  \draw[thick,fill=gray!15] (-0.5,-1) rectangle (0.5,1);
  \draw[thick,fill=orange!30] (0,0) circle (0.35);
  \draw[->,thick] (0.5,0) arc (0:300:0.5);
  \fill (0,0) circle (1.5pt);
  \node[align=center,font=\footnotesize,below] at (0,-1.15) {リアクションホイール\\重心が回転軸上：形が変わらない\\$\Gamma=$ 一定、曲率 $F=0$\\何回転でも回せる};
\end{scope}
\begin{scope}[shift={(5,0)}]
  \draw[thick,fill=gray!15] (-0.5,-1) rectangle (0.5,1);
  \fill (0,0.6) circle (2pt);
  \draw[line width=3pt,blue!60] (0,0.6)--++(20:1.3);
  \draw[line width=3pt,blue!25] (0,0.6)--++(80:1.1);
  \draw[->,thick] (0,0.6)++(20:0.9) arc (20:80:0.9);
  \node[align=center,font=\footnotesize,below] at (0,-1.15) {手足（AMBAC）\\重心が動く：形が変わる\\$\Gamma(q)$ が形に依存、$F\neq0$\\可動域がある};
\end{scope}
\node[font=\footnotesize] at (2.5,0.4) {$\dot\theta=-\Gamma\,\dot q$};
\end{tikzpicture}
